\section{Related Work}
\label{sec:related_work}
Planning charging infrastructure has been studied extensively, and recent reviews organize the field along the dimensions of station siting, sizing, grid reinforcement, trip-considerations and the treatment of uncertainty \cite{review2025planning, pagany2019review}.  Most contributions optimize either where demand is best served or how the grid copes with the additional load; bidirectional charging, in which vehicles also feed power back, is only beginning to be integrated into such planning frameworks \cite{v2gplanning2023}.

\subsection{Traffic-Aware Siting}

Traffic-aware siting methods estimate charging demand from origin-destination flows, activity patterns, and route choices, often placing stations so as to maximize coverage of traffic along routes. Data-driven analyses suggest that charging demand is largely shaped by proximity to amenities, EV density, and adjacency to high-traffic corridors \cite{causal2025placement}, which motivates deriving demand from land-use data as we do. Microscopic traffic simulators are increasingly used for this purpose: SUMO in particular provides built-in energy and charging-infrastructure models \cite{kurczveil2014sumo} and has been used to assign electric vehicles to charging stations from simulated traffic data \cite{sumoassignment2023}. Such approaches improve accessibility and reduce detours, but typically abstract away local constraints in the distribution grid.

\subsection{Grid-Aware Siting}
Grid-aware planning approaches prioritize feeder hosting capacity, voltage stability, and peak-load mitigation, frequently relying on metaheuristic optimization to jointly site and size stations under electrical constraints \cite{review2025planning}. While they improve technical feasibility, they can yield a poor user experience when travel behavior is not modeled in sufficient detail. The placement of charging stations is often motivated not only by charging demand and by economic advantages, but also by repercussions for the electricalnetwork  \cite{ahmad2022location}. While a reasonable amount of studies have been conducted on the siting of charging stations, there are only few that consider working with large-scale eletrical vehicle charging infrastructure \cite{largescale2024infrastructure}.

\subsection{Coupled and Bidirectional Approaches}
A smaller but growing body of work couples transportation simulation with power-system analysis, for example through co-simulation frameworks that pass simulated charging loads to a distribution-grid model \cite{cosim2025grid}. For bidirectional charging, planning models increasingly distinguish public and private chargers and account for uncertainty in demand and renewable generation \cite{v2gplanning2023}. Integrated and openly reproducible frameworks that combine microscopic mobility simulation with bidirectional charging at the level of individual public stations and private home wallboxes, however, remain limited. This paper addresses that gap on the mobility-and-demand side, while leaving a detailed electrical model of the distribution grid to future work.
